\section{Reivindicações -- minuta preliminar}

\textbf{Importante}: as reivindicações abaixo são rascunho técnico. Devem ser revisadas por profissional especializado antes do depósito.

\begin{enumerate}
  \item Método implementado por computador para priorização de causas de evasão escolar, caracterizado por compreender:
  \begin{enumerate}[label=\alph*)]
    \item armazenar, em memória computacional, dados de um projeto educacional, grupos participantes e causas possíveis de evasão;
    \item gerar, por processador, um primeiro conjunto de instruções para criação de formulário digital de avaliação das causas e comparação entre grupos participantes;
    \item receber arquivo de respostas do referido formulário digital;
    \item agregar, para cada causa possível, frequências de respostas em escala ordinal;
    \item calcular, para cada causa possível, uma mediana discreta a partir das frequências agregadas;
    \item calcular, para cada causa possível, uma concordância percentual a partir de notas previamente selecionadas como positivas;
    \item aprovar uma causa quando a mediana discreta atende a um limite mínimo e a concordância percentual atende a um limite mínimo;
    \item calcular pesos dos grupos participantes a partir de comparações pareadas;
    \item gerar um segundo conjunto de instruções para criação de formulário digital contendo as causas aprovadas;
    \item receber arquivo de respostas finais relativo às causas aprovadas; e
    \item gerar, com base nos pesos dos grupos participantes e nas respostas finais, ranking ordenado das causas aprovadas.
  \end{enumerate}

  \item Método, de acordo com a reivindicação 1, caracterizado pelo fato de que a escala ordinal compreende notas de 1 a 5.

  \item Método, de acordo com qualquer uma das reivindicações anteriores, caracterizado pelo fato de que as notas positivas são configuráveis pelo usuário.

  \item Método, de acordo com qualquer uma das reivindicações anteriores, caracterizado pelo fato de que a mediana mínima é configurável em valores inteiros de 1 a 5.

  \item Método, de acordo com qualquer uma das reivindicações anteriores, caracterizado pelo fato de que a concordância mínima é configurável em percentual inteiro.

  \item Método, de acordo com qualquer uma das reivindicações anteriores, caracterizado pelo fato de que os pesos dos grupos participantes são obtidos a partir de matriz recíproca formada por comparações pareadas.

  \item Método, de acordo com a reivindicação 6, caracterizado por calcular uma razão de consistência das comparações pareadas e sinalizar inconsistência quando a razão excede limite configurável.

  \item Método, de acordo com qualquer uma das reivindicações anteriores, caracterizado por gerar relatório contendo pelo menos: número de respondentes, causas aprovadas, critérios de aprovação, pesos dos grupos e ranking final.

  \item Sistema computacional para priorização de causas de evasão escolar, caracterizado por compreender:
  \begin{enumerate}[label=\alph*)]
    \item uma interface de cadastro de projetos, grupos participantes e causas possíveis;
    \item um módulo de geração de formulários digitais;
    \item um módulo de importação de respostas em CSV;
    \item um módulo de aprovação de causas por mediana e concordância;
    \item um módulo de cálculo de pesos de grupos participantes;
    \item um módulo de ranking final; e
    \item um módulo de relatório.
  \end{enumerate}

  \item Sistema, de acordo com a reivindicação 9, caracterizado pelo fato de que o módulo de geração de formulários digitais produz scripts executáveis em ambiente de criação de formulários online.

  \item Sistema, de acordo com a reivindicação 9 ou 10, caracterizado pelo fato de que o módulo de importação aceita arquivo CSV com nome arbitrário e armazena cópia interna padronizada no projeto.

  \item Meio de armazenamento legível por computador, não transitório, contendo instruções que, quando executadas por processador, realizam o método definido em qualquer uma das reivindicações 1 a 8.
\end{enumerate}
